% Source: https://artofproblemsolving.com/community/c3h2350400_proof_of_the_cheap_denesting_radicals_formula
% Stable identity: aops-topic-2350400-q01
\textbf{All the roots make sense.} Since $k^2=a^2-b<a^2$ and $a>0$, we have $0\le k<a$, so $\dfrac{a\pm k}{2}>0$; and $a^2-b\ge0$ gives $\sqrt b\le a$, so $a\pm\sqrt b\ge0$.

\textbf{Proof.} Put $P=\sqrt{\dfrac{a+k}{2}}$ and $Q=\sqrt{\dfrac{a-k}{2}}$. Then $P\ge Q\ge0$ (as $k\ge0$), $P^2+Q^2=a$ and
\[PQ=\sqrt{\dfrac{(a+k)(a-k)}{4}}=\dfrac{\sqrt{a^2-k^2}}{2}=\dfrac{\sqrt b}{2}.\]
Hence
\[(P\pm Q)^2=P^2+Q^2\pm2PQ=a\pm\sqrt b.\]
Both $P+Q$ and $P-Q$ are nonnegative, so they are the nonnegative square roots: $P+Q=\sqrt{a+\sqrt b}$ and $P-Q=\sqrt{a-\sqrt b}$. $\blacksquare$

\textbf{Why it is useful.} If $a$, $b$, $k$ are rational, the right-hand sides involve only square roots of the rational numbers $\dfrac{a\pm k}{2}$, so a ``nested'' radical $\sqrt{a\pm\sqrt b}$ is rewritten without nesting. (The condition is that $a^2-b$ is the square of a rational number; for other $a,b$ the formula is still true but $k$ itself is a square root and nothing is gained.) For example, with $a=3$, $b=8$ we get $k=1$ and $\sqrt{3+2\sqrt2}=\sqrt2+1$.
